%=============================================================
\chapter{Introdução}
\label{Intro}
%=============================================================

A pandemia de Covid-19, causada pelo vírus SARS-CoV-2, evidenciou a importância de prever a evolução de uma epidemia para apoiar as decisões de saúde pública. O vírus se espalha a partir de pequenas partículas líquidas emitidas por tosse, espirros e saliva de pessoas infectadas, e a sua disseminação está principalmente relacionada ao contato próximo de indivíduos em ambientes confinados, sem ventilação adequada e com interação frequente com um número significativo de pessoas. Pelas características do vírus, a evolução da doença se comporta de maneira diferente por região, e elementos específicos de cada área, como densidade populacional, pirâmide etária e disponibilidade de leitos hospitalares e de tratamento intensivo, são fundamentais para o modo como o espalhamento do vírus e a sua letalidade ocorrerão em tal região.

Para tentar prever como uma doença se espalhará em uma epidemia, são utilizados modelos epidemiológicos, que podem ser determinísticos, através de equações diferenciais, ou estocásticos, que evoluem através de probabilidades \citeonline{keeling2008}. Os modelos determinísticos clássicos descrevem a população de forma agregada e homogênea, e logo não capturam a heterogeneidade entre os indivíduos nem a estrutura espacial dos contatos. Neste trabalho será feita uma modelagem baseada em agentes, na qual a evolução da epidemia ocorre através da interação de indivíduos que possuem diversas características, como idade, sexo e outras que sejam relevantes ao modelo, distribuídos espacialmente em uma rede. Essa abordagem permite representar as particularidades de cada área e acompanhar não apenas o número de infectados, mas também a carga imposta ao sistema de saúde.

A presença de diversos parâmetros e o grande número de agentes para algumas áreas, que pode chegar à ordem de milhões de indivíduos, torna necessário um elevado poder computacional para a calibração e a implementação do modelo. Soma-se a isso o caráter estocástico do modelo, que exige a execução de um grande número de simulações independentes para que o comportamento médio da epidemia seja estimado com confiança. Como consequência, simular áreas grandes em uma implementação serial pode demandar muitas horas de processamento, e logo a paralelização do modelo torna-se uma forma eficiente de obter os seus resultados em tempo viável.

GPUs são orientadas a \textit{throughput} em detrimento de latência e, logo, conseguem realizar tarefas mais simples e em número maior do que CPUs. Como, no modelo baseado em agentes, todos os indivíduos da rede são processados de forma independente a cada passo de tempo, o problema é naturalmente adequado ao paralelismo de dados oferecido por essas placas. Com a popularização de GPUs Nvidia no mercado e a existência da API CUDA (\textit{Compute Unified Device Architecture}), que estende C/C++ para a computação em paralelo \cite{sanders2011}, a modelagem baseada em agentes em conjunto com a API CUDA torna-se uma boa alternativa para representar diferentes áreas do Brasil mantendo a eficiência, de modo que cada área não demande muito tempo de simulação.

Neste contexto, este trabalho apresenta a paralelização, utilizando CUDA, de um modelo baseado em agentes do tipo SEIR estendido para a Covid-19, aplicado a quatro áreas urbanas brasileiras com características distintas: São Paulo, Rocinha, Brasília e Manaus. A implementação paralela é desenvolvida a partir de uma implementação serial de referência, contra a qual é validada sob os mesmos parâmetros, e o seu ganho de desempenho é avaliado por meio da aceleração e da eficiência. Busca-se, assim, demonstrar que a paralelização preserva o comportamento do modelo e, ao mesmo tempo, reduz significativamente o tempo de simulação.

\section{Organização do trabalho}

Além desta introdução, este trabalho está organizado da seguinte forma. O Capítulo \ref{FT} apresenta os fundamentos teóricos necessários, abrangendo os modelos epidemiológicos compartimentais, a modelagem baseada em agentes, o método de Monte Carlo, a arquitetura de GPUs e o modelo de programação em CUDA. O Capítulo \ref{Met} descreve a metodologia, detalhando o modelo epidemiológico, os seus parâmetros e a estratégia de paralelização e validação. Os Capítulos \ref{cap:serial} e \ref{cap:gpu} apresentam, respectivamente, a implementação serial de referência e a implementação paralela em GPU. O Capítulo \ref{cap:resultados} apresenta e discute os resultados, validando a equivalência entre as implementações e avaliando o ganho de desempenho obtido.
